\bookchapter{第零章：嗨，开始之前}

\chapterintro{在正式探索前，咱们先聊点基本规则，确保后面说的你都懂。别急，磨刀不误砍柴工嘛。}

\booksection{一、这本书能帮你做什么？（和不能做的）}

它能帮你系统性地搞懂手机怎么用，明白功能背后的道理，而且无论你是安卓党还是苹果派，大多数思路是通用的。它会带着你动手练习，让你从“会用”升级到“懂用”。

它不是某一款特定手机的说明书，也不是追新品的科技新闻，更不是教你沉迷手机的指南。说实话，如果你期待的是“如何用手机刷视频更爽”，那可能要失望了——我们更关心的是如何让手机真正为你服务，而不是反过来。

\booksection{二、我们的“暗号”说明}

为了说得更清楚，书里会用一些小标记，这些标记就像我们之间的默契暗号：

\begin{itemize}
  \item \textbf{[按钮名]}：代表屏幕上具体的按钮或App名字，比如[设置]、[确定]。
  \item \textbf{>}：代表操作顺序，比如"点 [设置] > [声音]"，就是先点设置，再点声音。
  \item \textbf{加粗}：非常重要！可能是关键概念或安全警告，请务必留意。
  \item \emph{斜体}：一些补充的小想法或比喻，帮你更好地理解。
\end{itemize}

\booksection{三、安卓和iOS，咱都说说}

你猜怎么着，有些东西叫法不同，但干的事儿一样。比如：应用（安卓）和 App（iOS），设置（都叫设置），文件管理（安卓）和 文件（iOS）。知道是一回事就行，别被名词绕晕了。

另外，现在安卓阵营里还有华为的鸿蒙系统，它有自己的特色，后面我们会专门讲到。总之，不管你用什么手机，这本书都能帮到你。

\booksection{四、手机容量：为啥买了128G，可用却不到110G？}

容量单位大家知道：1TB = 1024GB，1GB = 1024MB，1MB = 1024KB……但这里有个坑：硬盘厂商通常按1000进制算，而系统按1024进制算，所以你买到的“128G”硬盘，实际容量会“缩水”。

但你手机容量缩水的原因主要有三个：

\begin{enumerate}
  \item \textbf{系统自己占了地方}：操作系统就像房子的承重墙，必须存在。现在的系统动辄占用10-20GB，这很正常。
  \item \textbf{厂商预装的应用}：新手机里自带的那些App也占空间，有的能卸载，有的不能。
  \item \textbf{分区和格式化损耗}：技术原因，就像书包的衬布也占体积。
\end{enumerate}

\begin{tcbinfo}{买手机的小建议}
现在128GB的手机基本不够用了，256GB起步比较稳妥。如果你喜欢拍照、玩游戏，建议直接上512GB甚至更大。记住，存储空间这东西，宁可富余，不可捉襟见肘。

\end{tcbinfo}

\begin{tcbpractice}{来动手：看看你的手机"身份证"}
打开 [设置]，找到关于手机（或类似选项），记下：手机型号、系统版本、总存储空间和剩余空间。这是了解它的第一步。顺便看看电池健康度，心里有个数。

\end{tcbpractice}

\booksection{五、你的手机是一座“数字资产库”}

在开始具体知识之前，我们先建立一个核心认知：\textbf{手机早已不只是通信工具，它是我们最重要的数字资产库}。联系人、聊天记录、家庭照片、工作文档、银行App、付款码、健康数据、电子证件……这些数据加在一起，价值往往超过手机本身。

\begin{tcbinfo}{为什么这很重要？}
如果手机丢了或坏了，你损失的不仅是几千块钱的硬件，\textbf{更可能是多年积累的不可再生数据}——孩子的成长照片、过世亲人的语音、没备份的重要合同。科学地整理与防护习惯，能从三个维度保护你：

\begin{itemize}
  \item \textbf{提升运行效率}：规范的文件分类让系统查找更快，手机长期不卡顿。
  \item \textbf{延长使用寿命}：合理的存储管理减少芯片读写负担，电池和存储都更耐用。
  \item \textbf{保障数据安全}：分层管理 + 定期备份 + 权限收敛，让"数据灾难"远离你。
\end{itemize}

\end{tcbinfo}

\booksubsection{我们的核心理念：“预防为主，防治结合”}

就像养生不能等到得病才看医生，\textbf{手机管理也要从“日常习惯”抓起}。本书反复强调的不是某个高级技巧，而是这些朴素的原则：

\begin{itemize}
  \item \textbf{预防为主}：装 App 只走官方商店、照片发之前先脱敏、密码不重复——这些"小动作"成本极低，但能挡住绝大多数风险。
  \item \textbf{定期体检}：每月一次整理下载文件夹、每季度一次权限审计、每半年一次全量备份——固定节奏比临时救火有效一百倍。
  \item \textbf{留有退路}：开启云同步、记好账号密码、写一份"数字遗产清单"——确保意外发生时你和家人都能从容应对。
\end{itemize}
